\section{Preliminaries}

Before detailing our ProtoDiff methodology, we first present the relevant background on few-shot classification, the prototypical network, and diffusion models.

\textbf{Few-shot classification.} We define the N-way K-shot classification problem, which consists of support sets $\mathcal{S}$ and a query set $\mathcal{Q}$. Each task $\mathcal{T}^i$, also known as an episode, represents a classification problem sampled from a task distribution $p(\mathcal{T})$. The \textit{way} of an episode denotes the number of classes within the support sets, while the \textit{shot} refers to the number of examples per class. Tasks are created from a dataset by randomly selecting a subset of classes, sampling points from these classes, and subsequently dividing the points into support and query sets. The episodic optimization approach~\cite{vinyals2016matching} trains the model iteratively, performing one episode update at a time.

\textbf{Prototypical network.} We develop our method based on the prototypical network (ProtoNet) by Snell \etal~\cite{snell2017prototypical}. Specifically, the ProtoNet leverages a non-parametric classifier that assigns a query point to the class having the nearest prototype in the learned embedding space. The prototype $\mathbf{z}^c$ of an object class $c$ is obtained by: $\mathbf{z}^c {=} \frac{1}{K}\sum_k f_{\phi}(\mathbf{x}^{c,k})$, where $f_{\phi}(\mathbf{x}^{c,k})$ is the feature embedding of the support sample $\mathbf{x}^{c,k}$, which is usually obtained by a convolutional neural network. For each query sample $\mathbf{x}^q$, the distribution over classes is calculated based on the softmax over distances to the prototypes of all classes in the embedding space:
\begin{equation}
\label{eqn:protonetProb}
    p(\mathbf{y}_{n}^q = c|\mathbf{x}^q) = \frac{\exp(-d(f_{\phi}(\mathbf{x}^q),\mathbf{z}^c))}{\sum_{c'} \exp(-d(f_{\phi}(\mathbf{x}^q),\mathbf{z}^{c'}))},
\end{equation}
where $\mathbf{y}^q$ denotes a random one-hot vector, with $\mathbf{y}_n^q$ indicating its $n$-th element, and $d(\cdot,\cdot)$ is some (Euclidean) distance function. Due to its non-parametric nature, the ProtoNet enjoys high flexibility and efficiency, achieving considerable success in few-shot learning.  To avoid confusion, we omit the superscript $c$ for the prototype $\mathbf{z}$ in this subsection.

\textbf{Diffusion model.}
In denoising diffusion probabilistic models~\cite{ho2020denoising}, a forward diffusion process, $q(\x_t|\x_{t-1})$, is characterized as a Markov chain that progressively introduces Gaussian noise at each time step $t$, beginning with a clean image $\x_0 \sim q(\x_0)$. The subsequent forward diffusion process is formulated as follows:
\begin{align}
\footnotesize
    q(\x_{T}|\x_0):=\prod_{t=1}^{T}q(\x_t|\x_{t-1}),\quad \mbox{where}\quad
    q(\x_t|\x_{t-1}):=\Nc(\x_t;\sqrt{1-\beta_t}\x_{t-1},\beta_t\Ib),
\end{align}
where $\{\beta\}_{t=0}^{T}$ is a variance schedule.
By defining $\alpha_t {:=} 1 {-} \beta_t$ and $\bar{\alpha_t} {:=} \prod_{s=1}^t\alpha_s$, the forward diffused sample at time step $t$, denoted as $\bm x_t$, can be generated in a single step as follows:
\begin{align}
    \label{eq:ddpm}
    \x_t=\sqrt{\bar{\alpha}_t}\x_0+\sqrt{1-\bar{\alpha}_t}\boldsymbol{\epsilon},\quad \mbox{where}\quad \boldsymbol{\epsilon}\sim\Nc(\textbf{0},\Ib).
\end{align}
Since the reverse of the forward step, $q(\x_{t-1}|\x_t)$, is computationally infeasible, the model learns to maximize the variational lower bound using parameterized Gaussian transitions, $p_\theta(\x_{t-1}|\x_t)$, where the parameter is denoted as $\theta$. Consequently, the reverse process is approximated as a Markov chain with the learned mean and fixed variance, starting from  a random noise $\x_T \sim  \Nc(\x_T; \textbf{0}, \boldsymbol{I})$:
\begin{align}
    p_\theta(\x_{0:T}):=p_\theta(\x_T)\prod_{t=1}^{T}p_\theta(\x_{t-1}|\x_{t}), 
\end{align}
\mbox{where}
\begin{align}\label{eq:mu}
     p_\theta(\x_{t-1}|\x_{t}):=\Nc(\x_{t-1};{\boldsymbol\mu}_\theta(\x_t,t),\sigma_t^2\Ib), \quad
    {\boldsymbol\mu}_\theta(\x_t,t):=\frac{1}{\sqrt{\alpha}_t}\Big{(}\x_t-\frac{1-\alpha_t}{\sqrt{1-\bar{\alpha}_t}}\boldsymbol{\epsilon}_\theta(\x_t,t)\Big{)}.
    %\quad     \boldsymbol{s}_\theta(\x_t,t)=-\frac{1}{\sqrt{1-\bar{\alpha}_t}}\z_\theta(\x_t,t)
\end{align}
Here, $\boldsymbol{\epsilon}_\theta(\x_t,t)$ is the diffusion model trained by optimizing the following objective function:
\begin{align}
    \label{eq:objective}
    \mathcal{L}_{\theta}=\mathbb{E}_{t,\x_0,\boldsymbol{\epsilon}}\Big{[}\|\boldsymbol{\epsilon}-\boldsymbol{\epsilon}_\theta(\sqrt{\bar{\alpha}_t}\x_0+\sqrt{1-\bar{\alpha}_t}\boldsymbol{\epsilon},t)\|^2\Big{]}.
\end{align}
Upon completing the optimization, the learned score function is integrated into the generative (or reverse) diffusion process. To sample from $p_\theta(\x_{t-1}|\x_t)$, one can perform the following:
\begin{align}\label{eq:reverse}
    \x_{t-1}=   {\boldsymbol\mu}_\theta(\x_t,t)+\sigma_t\boldsymbol{\epsilon}=\frac{1}{\sqrt{\alpha_t}}\Big{(}\x_t-\frac{1-\alpha_t}{\sqrt{1-\bar{\alpha}_t}}\boldsymbol{\epsilon}_\theta(\x_t,t)\Big{)}+\sigma_t\boldsymbol{\epsilon}.
\end{align}
%
In the case of conditional diffusion models~\cite{kwon2022diffusion, saharia2022palette, sasaki2021unit}, the diffusion model $\boldsymbol{\epsilon}_\theta$ in equations (\ref{eq:objective}) and (\ref{eq:reverse}) is substituted with $\boldsymbol{\epsilon}_\theta(\mathbf{c}, \sqrt{\bar{\alpha}_t}\x_0 + \sqrt{1 - \bar{\alpha}_t}\boldsymbol{\epsilon}, t)$, where $\mathbf{c}$ represents the corresponding conditions, e.g., other images, languages, and sounds, etc.
Consequently, the matched conditions strictly regulate the sample generation in a supervised manner, ensuring minimal changes to the image content.
Dhariwal and Nichol \cite{dhariwal2021diffusion} suggested classifier-guided image translation to generate the images of the specific classes along with a pre-trained classifier.
% Specifically, ${\boldsymbol\mu}\theta(\x_t,t)$ in equations (\ref{eq:mu}) and (\ref{eq:reverse}) is augmented with the gradient of the classifier: ${\bm \hat\mu}\theta(\x_t,t) := {\boldsymbol\mu}\theta(\x_t,t) + \sigma_t \nabla{\x_t} \log p_{\phi}(\y|\x_t)$.
Taking inspiration from conditional diffusion models that excel at generating specific images using additional information, we introduce the task-guided diffusion model that generates prototypes of specific classes by taking into account contextual information from various few-shot tasks. Note that our task-guided diffusion model operates not on an image $\x$, but on the prototype $\mathbf{z}$. 